\section{The wing constant, and existence}\label{sec:exist}

The log discrete local variance $\lambda=\log\dlv$ ranges over all of
$\R^{NM}$, so positivity is automatic. By Theorem~\ref{thm:pointwise} the model
is the stochastic differential equation
\begin{equation}\label{eq:abstract}
  \dd\lambda_t=\Bigl[-\ip{\eta_t}{b_t}-\tfrac12\abs{b_t}^{2}\Bigr]\dd t
   +b_t\,\dd W_t ,
\end{equation}
whose only non-elementary coefficient is $\eta=\psi/q$ from \eqref{eq:volsolve}
and \eqref{eq:eta}. Its denominator is small in the wings, and the question is
whether $\eta$ can blow up there.

The argument follows the route of Wissel's proof of his Theorem~3.5
\cite{Wissel2007}: log coordinates (his (3.29)), induction over expiries, and an
existence result for locally Lipschitz equations (his Proposition~3.7). Its one
non-standard ingredient is an a priori bound, independent of the state, on the
relative volatility of the calendar spread at the second expiry. Wissel could not
obtain this bound in general and imposes his condition (3.24) instead (his
Remark~1 following the proof, p.~18); in our coordinates the quantity is $\eta+b$
(Appendix~\ref{app:wissel}). We obtain it from an exact identity for the
volatility recursion and a known inequality for inverses of M-matrices.

\subsection{The volatility recursion}\label{sec:walk}\label{sec:master}

Write $P:=A_j(\dlv^j)^{-1}$, so that $q^j=Pq^{j-1}$.

\begin{corollary}[Positivity, unit column sums, absorbing ends]\label{cor:pos}
$P\ge0$ entrywise, $\one^{\!\top}P=\one^{\!\top}$, so $P_{im}\in[0,1]$, and the outer
columns of $P$ are $e_0$ and $e_{N+1}$.
\end{corollary}

\begin{proof}
By \eqref{eq:operator}, $A_j$ has non-positive off-diagonal entries, and in each
column the diagonal entry exceeds the sum of the moduli of the off-diagonal ones by
exactly one; so $A_j$ is a non-singular M-matrix and $P\ge0$
\cite{BermanPlemmons1994}. The column sums follow from $\one^{\!\top}A_j=\one^{\!\top}$
(Lemma~\ref{lem:null}), and the outer columns of $A_j$ are $e_0$ and $e_{N+1}$. The
first two statements are also Theorem~3.3 of \cite{BuehlerRyskin2016}, attributed
there to \cite{AndreasenHuge2011}.
\end{proof}

That is all that is needed about $P$ itself. What we want next is a recursion for
$\eta$ alone: one step per expiry, each adding a controlled amount, so that the
bound telescopes along the ladder. Dividing \eqref{eq:volsolve} by $q^j$ gives
\begin{equation}\label{eq:divided}
  \eta^{j}_i\;=\;\frac{(P\psi^{j-1})_i}{q^{j}_i}
   \;-\;\frac{\bigl(P\,G_j(q^j)\beta^j\bigr)_i}{q^{j}_i} ,
\end{equation}
the first term carrying $\eta^{j-1}$ and the second carrying $b^j$.

The first is an average. Writing $\psi^{j-1}=q^{j-1}\!\circ\eta^{j-1}$, its weights
$P_{ik}q^{j-1}_k/q^{j}_i$ are non-negative because $P\ge0$ and sum to one because
$q^j=Pq^{j-1}$, so it contributes at most $\lVert\eta^{j-1}\rVert_\infty$; they are
the matrix $\mathcal A$ of \eqref{eq:mastercomp}.

The source is the difficulty. Since $G_j$ is a second difference, row $i$ carries
the neighbouring densities $q^j_{i\pm1}$ as well as $q^j_i$, so dividing by $q^j_i$
leaves the ratios
\begin{equation}\label{eq:nbr}
  \frac{q^{j}_{i-1}}{q^{j}_i}
  \qquad\text{and}\qquad
  \frac{q^{j}_{i+1}}{q^{j}_i} ,
\end{equation}
which are unbounded in the wings, where the density decays fastest; nothing in
\eqref{eq:volsolve} bounds them. Let $E\in\R^{N\times(N+2)}$ restrict to the
interior nodes, so $EE^{\!\top}=I_N$, and recall
$A_j=I+\Gh_j\,\mathrm{diag}(\dlv^j)E$.

\begin{proposition}\label{prop:master}
With $b^j=\beta^j/\dlv^j$ and $\circ$ the entrywise product,
\begin{equation}\label{eq:master}
  \psi^{j}\;=\;P\,\psi^{j-1}\;-\;(I-P)\,E^{\!\top}\bigl(q^{j}\!\circ b^{j}\bigr),
\end{equation}
equivalently, with the entrywise non-negative
\[
  \mathcal A:=\mathrm{diag}(q^{j})^{-1}P\,\mathrm{diag}(q^{j-1}),
  \qquad
  \mathcal B:=\mathrm{diag}(q^{j})^{-1}P\,E^{\!\top}\mathrm{diag}(Eq^{j}),
\]
\begin{equation}\label{eq:mastercomp}
  \eta^{j}\;=\;\mathcal A\,\eta^{j-1}\;+\;\bigl(\mathcal B-E^{\!\top}\bigr)b^{j} ,
\end{equation}
where $\mathcal A\one=\one$, so the transport is an averaging operator, and the
$i$-th row of $\mathcal B$ sums to
$(Pq^{j})_i/q^{j}_i-\one_{\{i\in\{0,N+1\}\}}$.
\end{proposition}

\begin{proof}
Multiplying $A_j=I+\Gh_j\mathrm{diag}(\dlv^j)E$ by $P$ on the left and $E^{\!\top}$
on the right gives $P\Gh_j\mathrm{diag}(\dlv^j)=(I-P)E^{\!\top}$. By
\eqref{eq:volsolve} and \eqref{eq:factor},
$\psi^j=P\psi^{j-1}-P\Gh_j\mathrm{diag}(q^j)\beta^j$. Since $\beta^j=\dlv^j\!\circ b^j$
and diagonal matrices commute,
\[
  \mathrm{diag}(q^j)\beta^j
  =\mathrm{diag}(q^j)\mathrm{diag}(\dlv^j)\,b^j
  =\mathrm{diag}(\dlv^j)\mathrm{diag}(q^j)\,b^j
  =\mathrm{diag}(\dlv^j)\bigl(q^j\!\circ b^j\bigr),
\]
which is \eqref{eq:master}. Dividing it by $q^j_i$ gives \eqref{eq:mastercomp}: the rows of
$\mathcal A$ sum to $(Pq^{j-1})_i/q^j_i=1$, and for $\mathcal B$ the terms
$k\in\{0,N+1\}$ missing from $(Pq^j)_i$ vanish for interior $i$ and equal $q^j_i$
at $k=i$ for outer $i$, by Corollary~\ref{cor:pos}.
\end{proof}

So $\eta^j$ is an average of $\eta^{j-1}$ plus a source in $b^j$, and the march
across expiries is additive.

\begin{remark}[Why the source is harmless]\label{rem:increment}
The ratios \eqref{eq:nbr} disappear because the source is the operator's own
increment: by \eqref{eq:affine}, $G_j(q)\dlv=A_j(\dlv)q-q$. The mechanism is
clearest when $b^j\equiv c$, so that $\beta^j=c\,\dlv^j$:
\[
  G_j(q^j)\beta^j=c\bigl(A_jq^j-q^j\bigr)=c\bigl(q^{j-1}-q^j\bigr),
  \qquad
  P\,G_j(q^j)\beta^j=c\bigl(q^{j}-Pq^{j}\bigr),
\]
by $A_jq^j=q^{j-1}$ and $Pq^{j-1}=q^j$. Divided by $q^j_i$ this is
$c\bigl(1-(Pq^j)_i/q^j_i\bigr)$: the two ratios of \eqref{eq:nbr} have collapsed
to the single ratio $(Pq^j)_i/q^j_i$, whose size is what $\Theta_j$ measures in
\eqref{eq:rows}. For general $b$ the factorisation \eqref{eq:factor} moves $b$ off
$\dlv$ and onto $q$, diagonal matrices commuting, and the same computation applies
to $q^j\!\circ b^j$. This is motivation only: the proof above reaches
\eqref{eq:master} through $P\Gh_j\mathrm{diag}(\dlv^j)=(I-P)E^{\!\top}$ in one step.
\end{remark}

\subsection{The wing constant}\label{sec:bound}

That single surviving ratio, $(Pq^j)_i/q^j_i$, is what has to be bounded. Write
\begin{equation}\label{eq:rows}
  R^j_i:=\one_{\{1\le i\le N\}}+\frac{(Pq^j)_i}{q^j_i}-\one_{\{i\in\{0,N+1\}\}},
  \qquad
  \Theta_j:=\max_i R^j_i,
  \qquad
  \Xi:=\sum_{r=1}^{N}P_{rr}\le N ,
\end{equation}
and $H_j:=\max_i\abs{\eta^j_i}$. Here $R^j_i$ is the row bound of
\eqref{eq:mastercomp} at node $i$, and $\Theta_j$, the \emph{wing constant}, is its
size across the slab. The theorem shows $\Theta_j$ is bounded by the grid alone,
uniformly in the state; the bound rests on one inequality for inverses of
M-matrices and its consequence for $P^{2}$.

\begin{lemma}[Path product inequality]\label{lem:pathproduct}
For all nodes $i,r,m$,
\begin{equation}\label{eq:summand}
  P_{ir}\,P_{rm}\;\le\;P_{im}\,P_{rr}.
\end{equation}
\end{lemma}

\begin{proof}
$P$ is the inverse of an M-matrix by Corollary~\ref{cor:pos}, and every inverse
M-matrix satisfies the path product conditions
$a_{ij}a_{jk}\le a_{ik}a_{jj}$ for all $i,j,k$
\cite{ZhuZhangLiu2011}, a result first noted by Willoughby
\cite{Willoughby1977} and proved by Johnson and Smith \cite{JohnsonSmith1999}. Take
$(i,j,k)=(i,r,m)$.
\end{proof}

\begin{corollary}[Squared propagator]\label{cor:p2p}
For all $i,m$,
\begin{equation}\label{eq:p2p}
  \bigl(P^{2}\bigr)_{im}\le
  \begin{cases}
    \Xi\,P_{im}, & 1\le i\le N,\\
    (1+\Xi)\,P_{im}, & i\in\{0,N+1\}.
  \end{cases}
\end{equation}
\end{corollary}

\begin{proof}
Sum \eqref{eq:summand} over $r$. An outer $r$ contributes only if $r=i$, since
column $r$ of $P$ is $e_r$, and then contributes $P_{im}$.
\end{proof}

\begin{theorem}\label{thm:wing}
For every $\dlv\in(0,\infty)^N$,
\begin{equation}\label{eq:step}
  H_j\;\le\;H_{j-1}+\Theta_j\lVert b^j\rVert_\infty ,
\end{equation}
and
\begin{equation}\label{eq:thetabound}
  \Theta_j\le1+\Xi\le N+1,
  \qquad
  \Theta^{*}:=\sup_{\dlv}\sum_{j=2}^{M}\Theta_j\le(M-1)(N+1).
\end{equation}
\end{theorem}

\begin{proof}
Since $q^j=Pq^{j-1}\ge0$, Corollary~\ref{cor:p2p} gives
$(Pq^j)_i=\sum_m(P^2)_{im}q^{j-1}_m\le\Xi\,q^j_i$ for interior $i$ and
$\le(1+\Xi)q^j_i$ for outer $i$, so $R^j_i\le1+\Xi$ in both cases. Finally,
$\mathcal A\ge0$ with $\mathcal A\one=\one$ gives
$\lVert\mathcal A\eta^{j-1}\rVert_\infty\le H_{j-1}$, while row $i$ of
$(\mathcal B-E^{\!\top})b^{j}$ has modulus at most $\lVert b^j\rVert_\infty R^j_i$
by the row sums of Proposition~\ref{prop:master}.
\end{proof}

The bound is a property of the grid, not of $\dlv$, and it is far from tight on
realistic states: on Heston surfaces with $N=40$ the measured $\Theta_j$ is at most
$7.6$ against $41$.

\subsection{Existence}\label{sec:existence}

Iterating \eqref{eq:step} from the front slice,
\begin{equation}\label{eq:apriori}
  \sup_{j,i}\abs{\eta^j_{i,t}}
  \;\le\;\lVert\eta^1_t\rVert_\infty+(M-1)(N+1)\,\lVert b\rVert_\infty ,
\end{equation}
and iterating \eqref{eq:master} shows that $\eta^j$ is affine in $\eta^1$, with
coefficients depending on $\lambda$ only through $P_2,\dots,P_j$ and $q^2,\dots,q^j$.
The only $t$ on the right of \eqref{eq:apriori} is inside $\eta^1_t$, which enters
linearly, so it has to be integrable in time rather than bounded.

\begin{assumption}[Integrable front volatility]\label{ass:F}
For a horizon $T\le T_1$,
\begin{equation}\label{eq:F}
  \int_0^{T}\lVert\eta^{1}_t\rVert_\infty\,\dd t\;<\;\infty\qquad\text{almost surely.}
\end{equation}
\end{assumption}

\emph{TBD.}
